\chapter{Conversion disorder functional neurological disorder}
\index{Conversion disorder functional neurological disorder}

\medskip
\noindent\textit{Educational draft; seek professional advice for individual care.}

\section*{Definition and scope}
Conversion disorder functional neurological disorder describes a mental-health concern or behaviour that may cause distress, risk, or impaired functioning. Assessment should be respectful and non-stigmatising.

\section*{Contributors and risk}
Contributors may include stress, trauma, sleep disruption, physical illness, substance use, neurodevelopmental factors, isolation, and family history.

\section*{Presentation}
Assess mood, thoughts, behaviour, relationships, sleep, appetite, work or school function, substance use, and immediate safety.

\section*{Assessment}
Use history, mental-state examination, risk assessment, physical review when indicated, and validated questionnaires as aids.

\section*{Management}
Care may include psychoeducation, practical support, talking therapy, treatment of contributing conditions, medication, and a shared follow-up plan.

\section*{When to seek urgent care}
Seek urgent help for immediate danger, suicidal intent, psychosis, inability to meet basic needs, or rapid deterioration.

\section*{Prevention and follow-up}
Protect routines, maintain supportive contact, reduce harmful substance use, and arrange review when symptoms persist.

\section*{Sources and limitations}
This concise educational overview should be checked against current local clinical guidance and adapted to the individual.ection*{Safety-netting}
Seek urgent help for immediate danger, suicidal intent, psychosis, severe confusion, inability to meet basic needs, or rapidly worsening symptoms.
